% chapters/ch06.tex
\chapter{Validación del modelo y análisis de resultados}

\begin{refsection}
\addbibresource{references/ch06.bib}

\section{Propósito de la validación del modelo}

Una vez desarrollado el modelo de diagnóstico educativo inclusivo basado en datos y su correspondiente implementación tecnológica, resulta necesario evaluar su funcionamiento en contextos educativos reales. La validación del modelo tiene como objetivo analizar el grado en que el sistema tecnológico contribuye a la identificación de necesidades educativas diversas y al apoyo de la toma de decisiones pedagógicas fundamentadas.

En esta investigación, la validación del modelo se desarrolla mediante el análisis integrado de datos cuantitativos y cualitativos obtenidos durante la implementación del sistema en el contexto educativo seleccionado.

Este proceso permite evaluar:

\begin{itemize}
\item la capacidad del sistema para identificar patrones de aprendizaje relevantes;
\item la utilidad del diagnóstico generado para los docentes;
\item la coherencia entre los resultados del sistema y los principios del Diseño Universal para el Aprendizaje.
\end{itemize}

\section{Implementación del sistema en el contexto educativo}

El sistema tecnológico desarrollado fue implementado en un contexto educativo real con la participación de docentes y estudiantes seleccionados mediante muestreo intencional.

Durante la implementación se realizaron las siguientes actividades:

\begin{itemize}
\item aplicación de instrumentos diagnósticos digitales;
\item recopilación de datos educativos mediante el sistema tecnológico;
\item generación automática de perfiles de aprendizaje;
\item análisis de los resultados obtenidos por parte de los docentes.
\end{itemize}

La implementación permitió observar el funcionamiento del modelo en situaciones reales de enseñanza y aprendizaje, así como recoger información relevante para su evaluación.

\section{Resultados cuantitativos}

Los resultados cuantitativos se obtuvieron a partir del análisis de los datos recopilados por el sistema tecnológico durante el proceso de diagnóstico educativo.

Las variables analizadas incluyeron:

\begin{itemize}
\item patrones de participación en actividades educativas;
\item preferencias de aprendizaje;
\item indicadores de interacción con los recursos educativos;
\item resultados de los instrumentos diagnósticos.
\end{itemize}

\subsection{Identificación de patrones de aprendizaje}

El análisis de los datos permitió identificar diferentes patrones de aprendizaje entre los estudiantes participantes. Estos patrones reflejan la diversidad de estilos, ritmos y estrategias de aprendizaje presentes en el aula.

La Tabla \ref{tab:patrones-aprendizaje} presenta un ejemplo de los patrones identificados mediante el sistema.

\begin{table}[htbp]
\centering
\caption{Ejemplo de patrones de aprendizaje identificados mediante el sistema}
\label{tab:patrones-aprendizaje}

\begin{tabular}{lll}
\toprule
Perfil de aprendizaje & Características principales & Recomendaciones pedagógicas \\
\midrule
Perfil analítico & Preferencia por información estructurada & Uso de organizadores gráficos \\
Perfil visual & Alta respuesta a recursos visuales & Incorporación de material gráfico \\
Perfil colaborativo & Preferencia por trabajo en grupo & Actividades cooperativas \\
Perfil autónomo & Alta capacidad de autoaprendizaje & Proyectos individuales guiados \\
\bottomrule
\end{tabular}

\end{table}

Los resultados muestran que el sistema permite identificar patrones de aprendizaje relevantes que pueden ser utilizados por los docentes para ajustar sus estrategias pedagógicas.

\subsection{Apoyo a la toma de decisiones pedagógicas}

El análisis de los resultados generados por el sistema indica que la información proporcionada facilita la interpretación de las características de aprendizaje de los estudiantes.

En particular, el sistema permitió:

\begin{itemize}
\item identificar posibles barreras para el aprendizaje;
\item reconocer fortalezas educativas de los estudiantes;
\item sugerir estrategias pedagógicas alineadas con el \DUA.
\end{itemize}

Estos resultados respaldan la hipótesis de que los sistemas tecnológicos basados en datos pueden contribuir al fortalecimiento de la toma de decisiones pedagógicas.

\section{Resultados cualitativos}

Los resultados cualitativos se obtuvieron mediante entrevistas semiestructuradas realizadas a docentes participantes, así como mediante observación del uso del sistema en el aula.

El análisis de contenido permitió identificar varias categorías relevantes relacionadas con la percepción del sistema.

\subsection{Utilidad pedagógica del sistema}

Los docentes participantes señalaron que el sistema facilita la comprensión de las diferencias individuales entre los estudiantes.

En particular, destacaron que:

\begin{itemize}
\item el sistema proporciona información organizada sobre el aprendizaje de los estudiantes;
\item permite identificar necesidades educativas que no siempre son evidentes mediante observación directa;
\item facilita la planificación de estrategias pedagógicas diferenciadas.
\end{itemize}

\subsection{Percepción sobre el uso de datos en educación}

Otro aspecto relevante identificado en el análisis cualitativo se relaciona con la percepción de los docentes sobre el uso de datos educativos.

Los participantes señalaron que el uso de información basada en datos puede contribuir a mejorar la toma de decisiones pedagógicas, siempre que el sistema se utilice como herramienta de apoyo y no como sustituto del juicio profesional del docente.

Esta percepción coincide con planteamientos teóricos que destacan la importancia de integrar la analítica de datos con la reflexión pedagógica \parencite{selwyn2019}.

\section{Discusión de los resultados}

Los resultados obtenidos en la investigación permiten analizar el potencial del modelo de diagnóstico educativo inclusivo basado en datos para apoyar los procesos educativos.

En primer lugar, el sistema demostró ser capaz de identificar patrones de aprendizaje relevantes a partir del análisis de datos educativos. Este hallazgo respalda la hipótesis de que la analítica del aprendizaje puede contribuir a mejorar la comprensión de los procesos educativos.

En segundo lugar, los resultados sugieren que la información generada por el sistema puede apoyar a los docentes en la toma de decisiones pedagógicas orientadas a la inclusión educativa.

Finalmente, los resultados cualitativos indican que los docentes perciben el sistema como una herramienta útil para complementar sus prácticas pedagógicas, lo cual coincide con estudios recientes sobre el uso de tecnologías basadas en datos en educación \parencite{mandinach2016}.

\section{Limitaciones de la validación}

Como toda investigación aplicada, el proceso de validación del modelo presenta ciertas limitaciones que deben ser consideradas en la interpretación de los resultados.

Entre las principales limitaciones se encuentran:

\begin{itemize}
\item el tamaño limitado de la muestra;
\item el carácter contextualizado del estudio;
\item la naturaleza prototípica del sistema tecnológico desarrollado.
\end{itemize}

Estas limitaciones no invalidan los resultados obtenidos, pero sugieren la necesidad de futuras investigaciones que amplíen el alcance del modelo en otros contextos educativos.

\section{Implicaciones para la investigación y la práctica educativa}

Los resultados de la investigación tienen implicaciones tanto para el campo académico como para la práctica educativa.

En el ámbito académico, el estudio contribuye al desarrollo de modelos de diagnóstico educativo inclusivo basados en datos.

En el ámbito práctico, la investigación sugiere que los sistemas tecnológicos pueden desempeñar un papel relevante como herramientas de apoyo para la toma de decisiones pedagógicas orientadas a la inclusión educativa.

\printbibliography[heading=subbibliography,title={Referencias (Capítulo VI)}]

\end{refsection}